\label{sec:background}

\subsection{LODES Workplace Area Characteristics}

The Longitudinal Employer-Household Dynamics (LEHD) program at the U.S.\
Census Bureau publishes the Origin-Destination Employment Statistics (LODES)
dataset annually \citep{abowd2009lodes}.
LODES version~8 covers employment years 2002--2021 and provides, among
other files, the Workplace Area Characteristics (WAC) table.
The WAC file contains one row per census block, identified by a
15-character GEOID (\texttt{w\_geocode}), and reports the count of jobs
located in that block broken down by NAICS sector (\texttt{CNS01}--\texttt{CNS20})
as well as the total (\texttt{C000}).

LODES WAC is designed to characterize the economic activity \emph{located at}
a workplace, not the residences of the workers.
This distinction is critical for redistricting: we want to know what kind of
economic activity happens in a given tract, which is what WAC measures.
The complementary Origin Area Characteristics (OAC) file would measure the
economic activities of workers who live in each tract---a different, less
relevant question for geographic community character.

Block-level data are aggregated to census tracts by taking the first 11
characters of the 15-character block GEOID, which is the standard
Census Bureau encoding of the state (2 digits), county (3 digits),
and census tract (6 digits) identifiers.
After aggregation, each tract has a total job count and a breakdown across
20 NAICS sectors.

\paragraph{Data availability.}
WAC files are available for all 50 states from the LEHD portal
(\url{https://lehd.ces.census.gov/data/}).
The 2020 vintage used in this paper aggregates jobs reported during
calendar year 2020.
For a state like North Carolina, the WAC data covers 2,655 tracts after
block-to-tract aggregation; Wisconsin covers 1,527 tracts; Texas covers
6,877 tracts.
Tracts that appear in the census geography but have no WAC records are
purely residential (no workplace establishments) and receive a zero
character vector, handled gracefully by the implementation
(Section~\ref{sec:algorithm}).

\subsection{Related Work: Economic Geography and Redistricting}

The relationship between economic geography and partisan sorting has been
documented extensively in the redistricting literature.
\citet{rodden2019} shows that the residential segregation of urban and
rural economic communities---the dense, mixed-use employment centers of
cities versus the dispersed residential and agricultural economy of
exurbs---is a primary driver of the urban-suburban partisan gap that makes
geographic compactness criteria structurally disadvantageous to Democratic voters.
Standard geographic boundary-length weights encode and reinforce this
sorting, because the boundaries that METIS most avoids cutting are precisely
those that run along the dense urban-suburban economic transition.

Economic character weights offer a potential partial correction.
If commercial corridors and employment centers are used as cut seams rather
than as interior territory of residential districts, the resulting districts
may better reflect the actual community structure of the state.
This is not a partisan adjustment---no electoral data enters the
computation---but a community-structure adjustment that may incidentally
improve proportionality where economic and partisan geographies are aligned.

\citet{altman2011modeling} proposed using socioeconomic indicators as
community-of-interest proxies in redistricting.
Their approach relied on American Community Survey (ACS) socioeconomic
variables at the census-tract level.
LODES WAC provides a different and arguably more direct signal: observed
workplace employment by sector, rather than inferred socioeconomic status
from residential survey data.

\citet{stephanopoulos2012redistricting} argues for a ``territorial community''
framework in redistricting law that would require districts to reflect
coherent economic and social communities rather than merely minimizing
boundary disruption.
Economic character edge weights can be understood as a computational
implementation of this principle within the existing algorithmic framework.

\subsection{Edge Weight Design in Algorithmic Redistricting}

The \textsc{bisect} platform uses a three-layer compositor
\citep{bisect2024}.
Layer~2 (edge and vertex weights) has evolved over the course of the
B-track papers.
T.3 introduced subdivision-respecting (county) weights that strengthen
edges between tracts sharing the same county \citep{b10county}.
T.7 introduced VRA-aligned weights that strengthen edges between
minority-majority tracts to protect majority-minority districts \citep{b14vra}.
Economic character (M.9) introduces a third class of Layer~2 modification:
similarity-based weights that depend on the economic profile of the tracts
being connected rather than on their administrative membership or
demographic composition.

The cosine similarity formulation used here follows standard practice in
information retrieval and machine learning for comparing non-negative
document vectors \citep{manning2008ir}.
Its application to geographic tract comparison is, to our knowledge, novel
in the redistricting literature.
